\section{Related Work}

\subsection{AI Reviewer Profiles and Profile Depth Effects}

The use of large language models to simulate peer review has grown substantially as a strategy for scaling evaluation capacity in both academic and creative domains. Early work in this tradition optimizes for numeric accuracy relative to human judgments. Shah and Wein \cite{shah2023reviewers} investigate whether LLMs can replicate reviewer decisions on academic manuscripts, finding acceptable agreement on surface-level criteria but systematic degradation on assessments of novelty and contribution depth. Liu et al. \cite{liu2023reviewerbot} demonstrate that domain-targeted prompting improves the specificity of generated reviews, while the broader ``LLM-as-judge'' paradigm \cite{zheng2023judging} establishes that LLM evaluators can approximate human pairwise preferences under controlled conditions. Berry et al. \cite{berry2024aireviewer} extend this line to long-document structured review, showing that staged prompting improves review coherence but does not eliminate tendency toward generic praise in the absence of explicit evaluative schema.

Across this body of work, the target quantity is numeric proximity to human judgments, and persona is treated as an enabling mechanism rather than an independent variable. Profile richness is incidental rather than designed: reviewers are prompted with role descriptions sufficient to orient the model to the evaluation context, but the content and depth of those descriptions is not systematically varied. This leaves unexplored the question of whether profile specificity produces qualitatively different evaluative frameworks, not merely better-calibrated scores.

Our companion paper \cite{dellaLibera2025panelcreative} addresses this gap directly in the evaluation setting, establishing that reviewer profile depth produces a non-monotonic quality gradient: a sparse generic profile outperforms the bare baseline, but domain-specific profiles with named analytical frameworks produce evaluative language that does not appear at lower depth levels. Critically, that study finds that the C2 condition --- domain knowledge provided without evaluative framework --- scores below the C1 condition (task specification only), a result we term the domain-inversion effect. The present paper asks whether the same gradient structure obtains when the same profiles are applied to authoring rather than reviewing, and whether the C2 inversion replicates in the authoring direction.

\subsection{Persona Prompting for Creative Generation}

The application of persona prompting to creative tasks has been studied from several angles. Zhang et al. \cite{zhang2018personas} establish that explicit persona descriptions improve coherence and self-consistency in dialogue generation, confirming that LLMs respond to identity framing as more than surface decoration. Park et al. \cite{park2023generative} demonstrate that memory-augmented persona representations support persistent, coherent behavior across multi-turn generative agent simulations, suggesting that the benefits of persona extend beyond single-turn generation and depend on the continuity of identity context across the interaction. Dang et al. \cite{dang2023creative} apply persona conditioning to creative writing assistance, finding that distinct persona instantiations produce meaningfully differentiated stylistic outputs from the same underlying model, with implications for diversity in AI-assisted co-creative workflows.

The central assumption of this literature is that professional identity functions as a holistic unit: the researcher is either instantiated or not, and instantiation activates the model's relevant prior knowledge as a package. Quality variation within a persona is not measured, and the constituent components of professional identity --- domain knowledge, evaluative frameworks, philosophical disposition, characteristic cognitive moves --- are treated as a single bundle that is present or absent. Whether a richer, more internally differentiated persona description produces qualitatively better outputs than a thinner one is not addressed.

A second unexamined assumption is that reviewing expertise and authoring expertise are interchangeable along the persona axis. When a profile describes how an expert evaluates creative work, it is not self-evident that the same profile will serve as a useful generative guide. A reviewer brings a terminal orientation to finished artifacts; an author requires a generative orientation to work-in-progress. Whether reviewing expertise is generative or terminal --- whether it transfers cleanly to authoring, transfers partially, or introduces friction --- depends on properties of the profile's cognitive lens that prior work has not isolated or studied.

\subsection{The Authoring/Reviewing Distinction in Creativity Research}

Human creativity research provides a theoretical foundation for expecting the authoring/reviewing distinction to matter in practice. Csikszentmihalyi's systems model of creativity \cite{csikszentmihalyi1997creativity} distinguishes the creator (who generates novel variations), the domain (the symbol system in which work is expressed), and the field (the social group of experts who select which variations are preserved). Creator and field operate with structurally different knowledge: the creator requires procedural knowledge of how to produce within the domain's constraints, while the field applies declarative knowledge of what the domain values and why. A reviewer is, in this model, a field actor; an author is, at minimum, a domain actor whose procedural fluency is the primary requirement.

Weisberg's analysis of creative expertise \cite{weisberg2006creativity} reinforces this distinction at the cognitive level. Expert creative performance is not reducible to knowing what good work looks like; it requires internalized, practiced routines for generating candidate work and for diagnosing and repairing failures during generation. This is what Polanyi \cite{polanyi1966tacit} distinguishes as knowing-how from knowing-that: the tacit procedural knowledge that underlies skilled practice cannot be fully recovered from declarative statements about quality standards, however detailed those statements are.

The authoring/reviewing distinction maps directly onto this. A reviewing profile operationalizes knowing-that: it articulates, often with considerable precision, what properties a finished puzzle should have. An authoring profile, if it is to function generatively, must encode something closer to knowing-how: the procedures, heuristics, characteristic moves, and generative impulses that a skilled constructor brings to a blank page. Profile depth may improve performance in both the authoring and reviewing directions, but the profile's content --- specifically, whether its evaluative questions are \emph{generative} (invertible to authoring constraints) or \emph{terminal} (applicable only to finished artifacts) --- should determine how much reviewing expertise survives the transfer. We term this property the lens type of a profile, and we make it the primary explanatory variable in our designer comparison study.

\subsection{Puzzle Hunt Design as Creative Testbed}

Puzzle hunt design is a tractable testbed for empirical work on creative generation for the same reasons it is tractable for evaluation: each puzzle has a single verifiable extraction answer, so correctness can be confirmed independently of quality; the expert community has produced a documented body of evaluative philosophy; and domain data can be precisely controlled through world files that fix the available material before authoring begins.

Several properties make this testbed especially useful for studying profile depth effects. The domain is narrow enough to support precise expertise claims --- a reviewer with a stated specialization in puzzle construction either applies domain-appropriate evaluative criteria or does not, and this can be observed in the review text --- but rich enough that shallow expertise produces recognizably different output from deep expertise. The AoE2 corpus used in Paper~\cite{dellaLibera2025panelcreative} provides a fixed domain world against which both authoring and evaluation outputs can be assessed for world-model consistency, a criterion that would be unavailable in domains without closed corpora. Practitioner frameworks from Snyder \cite{snyder2019craftsmanship}, Katz \cite{katz2021crossword}, and Young \cite{young2018visual} provide the three designer profiles used in our comparison study, grounding the AI personas in authentic published perspectives rather than constructed approximations.

\medskip

Taken together, prior work establishes that AI reviewer personas improve evaluation quality, that persona prompting affects creative generation, and that the human creativity literature provides principled reasons to expect the authoring/reviewing distinction to matter. What no prior study has examined is whether the quality gradient observed for reviewer profiles extends symmetrically to authoring profiles, whether the domain-inversion effect replicates in the generative direction, and --- most distinctively --- which properties of a reviewing profile determine whether it transfers productively to authoring or introduces lens-type friction. These are the questions the present paper addresses.
